% 62e Olympiade Internationale de Mathématiques — Saint-Pétersbourg (épreuve à distance), Russie.
% Énoncés officiels en français (imo-official.org). Chaque problème vaut 7 points.

\begin{center}
  \emph{Premier jour --- lundi 19 juillet 2021 --- Durée : 4 heures 30}
\end{center}

\subsection*{Problème 1}
Soit $n \geqslant 100$ un entier. Clara écrit les nombres $n, n+1, \ldots, 2n$ sur des cartes
distinctes. Elle mélange ensuite les $n+1$ cartes, puis les sépare en deux piles. Montrer qu'au moins
l'une des piles contient deux cartes pour lesquelles la somme des nombres est un carré parfait.

\subsection*{Problème 2}
Montrer que l'inégalité suivante est vraie
\[
\sum_{i=1}^{n} \sum_{j=1}^{n} \sqrt{|x_i - x_j|} \leqslant \sum_{i=1}^{n} \sum_{j=1}^{n} \sqrt{|x_i + x_j|}
\]
pour tous les nombres réels $x_1, \ldots, x_n$.

\subsection*{Problème 3}
Soit $D$ un point à l'intérieur d'un triangle acutangle $ABC$ tel que $AB > AC$, vérifiant
$\widehat{DAB} = \widehat{CAD}$. Le point $E$ sur le segment $[AC]$ vérifie $\widehat{ADE} = \widehat{BCD}$, le point $F$
sur le segment $[AB]$ vérifie $\widehat{FDA} = \widehat{DBC}$ et le point $X$ sur la droite $(AC)$ vérifie $CX = BX$.
Soit $O_1$ et $O_2$ les centres des cercles circonscrits aux triangles $ADC$ et $EXD$ respectivement.
Montrer que les droites $(BC)$, $(EF)$ et $(O_1O_2)$ sont concourantes.

\bigskip
\begin{center}
  \emph{Second jour --- mardi 20 juillet 2021 --- Durée : 4 heures 30}
\end{center}

\subsection*{Problème 4}
Soit $\Gamma$ un cercle de centre $I$, et $ABCD$ un quadrilatère convexe tel que chacun des segments
$[AB]$, $[BC]$, $[CD]$ et $[DA]$ soit tangent à $\Gamma$. Soit $\Omega$ le cercle circonscrit au triangle $AIC$.
L'extension de $[BA]$ au-delà de $A$ coupe $\Omega$ en $X$, et l'extension de $[BC]$ au-delà de $C$ coupe $\Omega$ en $Z$.
Les extensions de $[AD]$ et $[CD]$ au-delà de $D$ coupent $\Omega$ en $Y$ et $T$ respectivement. Montrer que
\[
AD + DT + TX + XA = CD + DY + YZ + ZC.
\]

\subsection*{Problème 5}
Deux écureuils, Bushy et Jumpy, ont ramassé $2021$ noix pour l'hiver. Jumpy numérote les noix de $1$ à $2021$,
et creuse $2021$ petits trous disposés en cercle sur le sol autour de leur arbre favori. Le matin suivant,
Jumpy remarque que Bushy a placé une noix dans chaque trou, mais n'a pas fait attention à leur numérotation.
Insatisfait, Jumpy décide de réarranger les noix en effectuant une suite de $2021$ modifications. Lors de la
$k$-ième modification, Jumpy échange les deux noix adjacentes à la noix $k$. Montrer qu'il existe une valeur
de $k$ telle que, lors de la $k$-ième modification, Jumpy a échangé deux noix $a$ et $b$ telles que $a < k < b$.

\subsection*{Problème 6}
Soit $m \geqslant 2$ un entier, $A$ un ensemble fini d'entiers (non nécessairement positifs), et
$B_1, B_2, B_3, \ldots, B_m$ des sous-ensembles de $A$. On suppose que pour tout $k = 1, 2, \ldots, m$ la somme
des éléments de $B_k$ vaut $m^k$. Montrer que $A$ contient au moins $m/2$ éléments.
